\section{Interpretation and next experiment}

For this model pair and training budget, retaining the question as text is the
better-supported practical choice.
The translation pipeline is substantially more accurate than the latent bridge
under the numeric diagnostic, and a much smaller learned textual prefix performs
best on the controlled single-step tasks.
The bridge's synthetic source-shuffling effect and small-set memorization show
that a learned numerical connection is possible.
They do not establish that this interface transfers the translation model's
multilingual competence or improves the target's reasoning.

Several differences limit comparison with LangBridge's reported gains.
The source here is a causal translation decoder rather than a bidirectional
multilingual encoder; its intermediate states are compressed into fixed bins.
Alignment uses only 1,024 Japanese question--answer pairs with short numeric
targets, rather than a large English alignment corpus.
The source states are used at their native scales.
A post hoc descriptive calculation on the pooled training features gives RMS
values of 11.24, 16.95 and 0.726 for layers 4, 8 and 16, respectively.
Layer ablations therefore change both the represented information and its scale.
These observations identify differences to investigate, not demonstrated causes
of the accuracy gap.

The main limitations are the small generated training distribution, a single
source/target model pair, one training seed for each ablation and the exploratory
test exposure described in Section~\ref{sec:experiments}.
Bootstrap intervals describe item variation conditional on those fits, not broad
model or language generalization.
MGSM also differs substantially from the generated tasks, and pretraining
exposure to benchmark material cannot be ruled out.
The output cap and the last-number diagnostic further restrict interpretation of
the frozen models' long explanations.
Conversely, the learned textual prefix's strong single-step scores do not extend
to two-step questions or establish a natural-benchmark improvement.

A useful next experiment is to test whether the same bridge can reconstruct the
paired English question from its source states before fitting answer generation.
This would provide a direct check of information recovery through the interface
and denser alignment supervision.
It should be compared with direct answer fitting at a controlled data and compute
budget, while preserving an untouched evaluation set.
Normalizing native feature scales is a separate, inexpensive control before
attributing differences to source-layer depth.
Neither modification has been evaluated in the present study.

\section{Conclusion}

A frozen PLaMo translation model can be connected differentiably to a frozen
Qwen model through a small learned projection, and the connection can memorize
a tiny source-dependent answer mapping.
The tested affine, pooled interface nevertheless has low accuracy on both its
larger training set and held-out questions.
Ordinary text input, with translation or a learned textual prefix, is markedly
more effective in this setting.
The experiment supports further interface-alignment diagnostics rather than a
claim of improved multilingual reasoning from hidden-state injection alone.
